\chapter{2022:James Maynard}\label{ch2022jamesmaynard:james-maynard-fields-medal-2022-a-comprehensive-report-on-his-work-and-impact}

\begin{center}\rule{0.5\linewidth}{0.5pt}\end{center}

\textbf{Prepared:} July 2026

\textbf{Subject:} James Maynard (b.\ June 10, 1987, Essex, England)

\textbf{Position:} Professor, University of Oxford (Fellow of St Catherine's College)

\textbf{Doctoral Advisor:} Roger Heath-Brown (University of Oxford)

\textbf{Fields Medal Citation:} \emph{``For contributions to analytic number theory, which have led to major advances on the structure of numbers in arithmetic progressions and in the understanding of their arithmetic structure.''}

\begin{center}\rule{0.5\linewidth}{0.5pt}\end{center}

\begin{prereqs}
Analytic number theory (sieve methods, exponential sums, Dirichlet series), basic algebraic number theory, elementary probability, complex analysis.
\end{prereqs}

\begin{objectives}
Explain the bounded gaps result and subsequent improvements; describe the work on large gaps between primes; understand the method of densities and its applications to equidistribution of primes.
\end{objectives}

\section{Problem History}

The prime number theorem states that the number of primes up to $x$ is $\pi(x) \sim x/\log x$. The twin prime conjecture asks: are there infinitely many primes $p$ such that $p+2$ is also prime? More generally, for any even integer $2k$, are there infinitely many prime pairs $(p, p+2k)$? This question dates back to at least the 19th century and remains one of the most famous open problems in number theory. The conjecture has resisted all attempts at proof despite its deceptively simple statement.

The first major breakthrough was Zhang's 2014 result: there exists a constant $H$ such that there are infinitely many intervals $[x, x+H]$ containing at least two primes. Zhang showed $H = 70,000,000$. The proof uses the GPY (Goldston--Pintz--Y{\i}ld{\i}r{\i}m) method combined with the Bombieri--Vinogradov theorem for smooth moduli. Before Zhang's work, it was not even known whether bounded gaps between primes existed for any finite bound, making this a watershed moment in analytic number theory.

The GPY sieve requires a strong level of distribution for primes in arithmetic progressions, which was not available for all moduli. Zhang circumvented this by using the Bombieri--Vinogradov theorem for smooth moduli rather than for all moduli. The question of \textbf{large gaps between primes}---how large $p_{n+1} - p_n$ can be---was also open. Cram\'er's conjecture (1937) predicts $p_{n+1} - p_n = O((\log p_n)^2)$, but it was not even known whether $\limsup (p_{n+1} - p_n) / \log p_n > 1$ unconditionally.

The \textbf{Elliott--Halberstam conjecture} (1968) predicts that the primes have a strong level of distribution for all moduli, which would imply bounded gaps with $H \le 12$. This conjecture, if true, would dramatically improve all known bounds. The Bombieri--Vinogradov theorem provides a level of distribution of $1/2$, but the Elliott--Halberstam conjecture asserts a level of distribution of $1$. Proving this remains one of the central challenges in analytic number theory.

The \textbf{Szemer\'edi--Tur\'an theorem} in combinatorics states that any subset of the integers with positive upper density contains arbitrarily long arithmetic progressions. The primes have zero density, so this theorem does not directly apply. However, Green and Tao (2004) proved that the primes contain arbitrarily long arithmetic progressions by developing a transference principle that allows combinatorial results to be applied to sparse sets. This transference principle, and the related concept of pseudorandom measures, play an important role in Maynard's subsequent work on bounded gaps.

The \textbf{Dickson conjecture} (1904) generalizes the twin prime conjecture: for any admissible tuple $(h_1, \ldots, h_k)$ of integers, there are infinitely many integers $n$ such that $n + h_1, \ldots, n + h_k$ are all prime. The notion of admissibility requires that for every prime $p$, the set $\{h_1 \bmod p, \ldots, h_k \bmod p\}$ does not cover all residue classes modulo $p$. The twin prime conjecture is the case $k = 2$, $h = (0, 2)$. Proving the Dickson conjecture in full generality remains beyond the reach of current methods, but partial results have been obtained.

The \textbf{Bombieri--Vinogradov theorem} (1965) is a strengthening of the prime number theorem in arithmetic progressions: for any $A > 0$, we have $\sum_{q \le x^{1/2}/(\log x)^A} \max_{(a,q)=1} |\pi(x;q,a) - \operatorname{li}(x)/\phi(q)| = o(x/\log x)$. This theorem controls the error term in the distribution of primes in arithmetic progressions on average over moduli and is essential for all modern sieve methods.

The fundamental difficulty in proving bounded gaps is that the primes, unlike sequences with positive density, are too sparse for the Szemer\'edi--Tur\'an theorem to apply directly. The primes have density $1/\log x$ near $x$, which tends to zero. The sieve methods of Brun, Selberg, and Goldston--Pintz--Y{\i}ld{\i}r{\i}m are designed precisely to handle such sparse sequences, but they require analytic inputs about the distribution of primes in arithmetic progressions that are difficult to obtain unconditionally.

The \textbf{level of distribution} of a sequence is a measure of how well the sequence is distributed among arithmetic progressions. For the primes, the Bombieri--Vinogradov theorem gives a level of distribution of $1/2$, meaning the primes are well-distributed among progressions of moduli up to $x^{1/2}$. The Elliott--Halberstam conjecture would give a level of distribution of $1$. The level of distribution directly determines how many primes can be detected by sieve methods, and hence how small the gaps between primes can be shown to be.

The historical development of sieve methods began with Brun's sieve (1919), which showed that the sum of reciprocals of twin primes converges (Brun's constant). Selberg's upper bound sieve (1950) provided a simpler alternative to the sieve of Eratosthenes--Legendre and gave the best upper bounds for many problems. The GPY method (2005) combined these sieve techniques with the Bombieri--Vinogradov theorem to show that $\liminf (p_{n+1} - p_n) / \log p_n \le 16$ under the Elliott--Halberstam conjecture. This was the starting point for Zhang's breakthrough and Maynard's subsequent improvement.

The connection between sieve methods and the theory of \textbf{exponential sums} is essential. The distribution of primes in arithmetic progressions is controlled by character sums $\sum_{n \le x} \chi(n)$, where $\chi$ is a Dirichlet character modulo $q$. The Polya--Vinogradov inequality gives the bound $|\sum_{n \le x} \chi(n)| \ll \sqrt{q} \log q$ for primitive characters. Stronger bounds, known as Burgess-type estimates, give $|\sum_{n \le x} \chi(n)| \ll x^{1-1/r} q^{(r+1)/(2r^2)}$ for any $r \ge 1$. These estimates are used to control the error terms in the sieve and are the reason why the level of distribution of primes is so important.

\section{How Was It Solved}

\textbf{Improved gaps between primes (2013).} Maynard improved the GPY method by replacing the requirement of primes in short intervals with a density condition on primes in arithmetic progressions. He introduced a multidimensional sieve using positive-definite quadratic forms. For $m \ge 1$, the sieve selects $m$ primes from $2m$ linear forms $L_i(n) = q_i n + b_i$. The key innovation is optimizing over the choice of quadratic form rather than fixing it, which allows for a significant improvement in the density requirement. The quadratic form determines how the $2m$ linear forms are weighted, and the optimization over this choice provides more efficient weights than the original GPY construction.

\textbf{Bounded gaps between primes (2014).} Maynard proved that for any $m \ge 1$, there are infinitely many intervals $[x, x + H(m)]$ containing at least $m$ primes, with $H(m) \le C m^3 2^{2m}$. This implies $H(2) \le 246$ for infinitely many prime pairs, a dramatic improvement over Zhang's original bound of $70,000,000$. The proof proceeds by constructing a multidimensional sieve and optimizing the weights using the theory of positive-definite quadratic forms. The key estimate is a lower bound for the integral of the sieve weight over the diagonal, which requires understanding the distribution of primes in arithmetic progressions with large moduli.

The improvement over GPY comes from several sources. First, the multidimensional sieve allows one to detect clusters of primes rather than individual primes. Second, the optimization over quadratic forms provides more efficient weights than the GPY weights. Third, the density condition is weaker than the Bombieri--Vinogradov theorem, allowing the result to be proved unconditionally. The conceptual leap was to observe that by using a multidimensional sieve with $m$ dimensions, one could detect clusters of $m$ primes simultaneously, which is a weaker requirement than detecting individual prime gaps.

The technical core of Maynard's argument is the optimization of a positive-definite quadratic form $Q$ on $\mathbb{R}^{2m}$ subject to a constraint arising from the level of distribution of primes. The sieve weight is defined using the quadratic form as $w(n) = \left( \sum_{i=1}^{2m} r_i(n) \right)^m$, where $r_i(n)$ is the indicator function that $L_i(n)$ is prime. The optimization over $Q$ minimizes the contribution of the main term while keeping the diagonal contribution large enough to produce a positive lower bound for the number of prime clusters.

\textbf{Large gaps between primes (2017).} With Ford, Green, Konyagin, and Tao, Maynard proved $\limsup (p_{n+1} - p_n) / (\log p_n \cdot (\log\log p_n)^2 \cdot \log\log\log\log p_n / (\log\log\log p_n)^2) \ge e^\gamma$. This dramatically improved the previous lower bounds of Rankin (1938) and others, which had additional logarithmic factors in the denominator. The proof uses a carefully constructed set of integers with few prime factors and estimates for the distribution of smooth numbers.

The construction of the prime-free interval uses the \textbf{sieve of Eratosthenes} in a refined form: one shows that every integer in $[x, x+H]$ has a prime factor from a specific set $S$ of primes, where $S$ is chosen so that $\sum_{p \in S} 1/p$ is large but the sieving is efficient. The estimate for the number of integers remaining after sieving uses the theory of \textbf{smooth numbers}---integers with no prime factors exceeding a given bound---and the distribution of primes in arithmetic progressions.

The key technical tool in the large gaps result is the \textbf{smooth number estimate} of Hildebrand and Tenenbaum: the number of $y$-smooth integers up to $x$ is $\Psi(x, y) = x \rho(u)$, where $u = \log x / \log y$ and $\rho$ is the Dickman function. This estimate is used to show that a carefully chosen set of integers has many smooth numbers, which in turn implies the existence of long prime-free intervals. The Dickman function satisfies the differential-difference equation $u \rho'(u) = -\rho(u-1)$ for $u > 1$, with $\rho(u) = 1$ for $0 \le u \le 1$.

\textbf{The discrete logarithm conjecture (2019).} Maynard proved the discrete logarithm conjecture for short intervals, showing that the number of elements in a short interval of the multiplicative group of a finite field is close to the expected value. This result has applications to coding theory and to the distribution of elements in finite fields.

\textbf{Multiples of a given integer (2022).} Maynard proved that if $N$ is any fixed positive integer, almost all intervals $[X, X + N]$ contain a prime. This resolves a conjecture of Erd\H{o}s for the case $N = 1$. For general $N$, the result provides a quantitative bound on the number of intervals that do not contain a prime.

The result for multiples is particularly interesting because it connects to the broader question of how primes are distributed among residue classes. For a fixed $N$, the number of intervals $[X, X+N]$ that do not contain a prime is related to the density of integers that are divisible by all primes up to $N$. The proof uses a refinement of the Maynard sieve adapted to the specific structure of short intervals.

The connection to the \textbf{Legendre conjecture}---that there is always a prime between $n^2$ and $(n+1)^2$---is worth noting. While Maynard's methods do not resolve this conjecture (which would require proving bounded gaps between primes in a much stronger sense), the density method provides tools for studying the distribution of primes in short intervals of the form $[n^2, (n+1)^2]$. The best known result in this direction is due to Ingham (1937), who showed that there is a prime between $n^3$ and $(n+1)^3$.

\textbf{The density method.} Maynard's density method is a general framework for proving that sets of integers contain arithmetic structures. Given a set $A$ of integers with positive density, one wishes to show that $A$ contains $m$-term arithmetic progressions for all $m$. The method uses a weight function adapted to $A$ and estimates the main term and error terms separately. The main term is positive if $A$ has positive density, and the error terms are controlled using exponential sum estimates.

The density method was developed in collaboration with Tao and builds on the transference principle of Green and Tao. The key insight is that one can transfer results from dense sets to sparse sets by constructing a ``pseudorandom majorant'' for the sparse set. The pseudorandom majorant is a function that majorizes the indicator function of the sparse set and satisfies strong mixing conditions. These mixing conditions allow one to use the ergodic-theoretic machinery developed for dense sets.

\section{Impact of This Discovery}

The bounded gaps result is one of the most celebrated achievements of 21st-century number theory. The multidimensional Maynard sieve has become a standard tool in analytic number theory, used extensively in problems involving distribution of primes in arithmetic progressions. The result answered a question that had been open since the 19th century and demonstrated that the distribution of primes is more regular than previously expected.

Maynard's work on large gaps between primes provides the best known lower bounds, though the conjectured quadratic gap remains unproved. The techniques used to construct prime-free intervals have connections to the theory of smooth numbers and have been applied to other problems in additive combinatorics and additive number theory.

The Maynard--Tao method has been extended to number fields by Shao, establishing bounded gaps between primes in the Chebotarev density setting. This extension requires new techniques for handling the splitting of primes in number fields and the distribution of Frobenius elements.

The impact on \textbf{computational number theory} has been significant. The explicit bound $H(2) \le 246$ provides a concrete target for computational verification. Computational evidence strongly supports the conjecture that $H(2) = 2$ (the twin prime conjecture), and the search for prime constellations of length $2$ within intervals of length $246$ has been carried out extensively. The first occurrence of a gap of size $2$ or less beyond $10^{18}$ has been verified computationally, and the distribution of small gaps appears to follow the Hardy--Littlewood conjecture.

The \textbf{Selberg sieve} (1950) is the starting point for understanding Maynard's method. The Selberg sieve provides an upper bound for the number of primes in an arithmetic progression by optimizing a quadratic form over a set of parameters. Maynard's innovation was to observe that this optimization can be performed in multiple dimensions, which provides more efficient weights and leads to better bounds on prime gaps. The connection between the Selberg sieve and the multidimensional sieve was not apparent before Maynard's work.

The density method has applications beyond prime number theory, including the study of Diophantine equations, the distribution of values of arithmetic functions, and the theory of additive combinatorics. The method provides a unified framework for studying the distribution of integers satisfying local conditions.

The \textbf{primes in short intervals} problem asks for bounds on the number of primes in $[x, x + x^\theta]$ for small $\theta > 0$. The prime number theorem implies that $\pi(x + x^\theta) - \pi(x) \sim x^\theta / \log x$ for any $\theta > 0$, but the best unconditional result for $\theta$ close to $0$ is due to Huxley (1972), who showed the result holds for $\theta > 7/12$. Maynard's work on bounded gaps is complementary: it shows that there are infinitely many intervals of bounded length containing at least $2$ primes, but does not address whether most intervals of a given length contain primes.

The \textbf{Landau's problems} ask four classical questions about primes: (1) are there infinitely many primes of the form $n^2 + 1$? (2) are there infinitely many primes of the form $n^2 + m^2 + 1$? (3) are there infinitely many Sophie Germain primes (primes $p$ such that $2p+1$ is also prime)? (4) are there infinitely many primes of the form $n^2 + m^2$? While Maynard's methods do not directly resolve these problems, they provide the sieve-theoretic framework within which progress on these conjectures is made. The Brun sieve gives an upper bound of $O(x^{1/2} / \log x)$ for the number of primes of the form $n^2 + 1$ up to $x$, but the conjectured answer is $c \sqrt{x} / \log x$ for an explicit constant $c$.

The \textbf{Parisi conjecture} in probability theory and the \textbf{Sato--Tate conjecture} in number theory are related to the distribution of Fourier coefficients of modular forms, which are intimately connected to the distribution of primes. Maynard's work on bounded gaps provides evidence for the strong regularity of prime distribution that underlies these conjectures.

\section{Modern Views and Improvements}

Maynard's method continues to yield new results. Extensions to more general number fields and to higher-dimensional varieties are active areas of current research. The Maynard--Tao method has been adapted to study primes represented by binary quadratic forms and primes in orbits of dynamical systems.

The relationship between the bounded gaps problem and the Elliott--Halberstam conjecture remains an active area: proving the full Elliott--Halberstam conjecture would give bounded gaps with $H \le 12$. Recent work of Frei and Maynard has improved the level of distribution in certain settings, bringing the bound closer to the conjectured value.

Maynard's work on the distribution of primes near squares, and on the gaps between almost-primes (integers with a bounded number of prime factors), continues to produce new results. The gaps between $k$-almost-primes are significantly smaller than the gaps between primes, and Maynard's methods provide the best known bounds for these gaps. The interaction with the theory of the Liouville function and the Chowla conjecture represents another frontier of current research.

The \textbf{Chowla conjecture} predicts that the Liouville function $\lambda(n) = (-1)^{\Omega(n)}$ (where $\Omega(n)$ is the number of prime factors of $n$ counted with multiplicity) has no correlations: $\sum_{n \le x} \lambda(n+h_1) \cdots \lambda(n+h_k) = o(x)$ for any $k$ and any distinct $h_1, \ldots, h_k$. This conjecture, if true, would imply the twin prime conjecture and much more. Maynard and Tao have made partial progress on this conjecture using methods related to the density increment argument.

The \textbf{Goldston--Pintz--Y{\i}ld{\i}r{\i}m (GPY) method}, which preceded Maynard's work, used a one-dimensional sieve to show that $\liminf (p_{n+1} - p_n) / \log p_n = 0$ under the Elliott--Halberstam conjecture. Maynard's multidimensional improvement showed the same conclusion unconditionally. The conceptual leap was to observe that by using a multidimensional sieve, one could detect clusters of primes simultaneously, which is a weaker requirement than detecting individual prime gaps.

The \textbf{distribution of gaps between consecutive primes} is closely related to Maynard's work. For a random sequence of primes with density $1/\log x$, the gaps would follow a Poisson distribution with mean $\log x$. The actual distribution of gaps between consecutive primes shows deviations from this model, particularly for small gaps. Maynard's result implies that the distribution function $F(t) = \lim_{x \to \infty} \#\{n \le x : p_{n+1} - p_n \le t \log p_n\} / \#\{n \le x\}$ satisfies $F(0) > 0$, which is impossible for a Poisson distribution. This deviation is one of the deepest manifestations of the structured nature of the primes.

The interaction between Maynard's sieve methods and the theory of \textbf{automorphic forms} is an emerging area. The distribution of Fourier coefficients of modular forms, and their relationship to primes in arithmetic progressions, provides a natural setting for applying Maynard's methods. The subconvexity bounds for $L$-functions provide the analytic input needed to control the error terms in the sieve.

The \textbf{quantitative distribution of primes in tuples} is an active area that builds directly on Maynard's work. The Hardy--Littlewood $k$-tuple conjecture predicts the precise asymptotic for the number of integers $n \le x$ such that $n + h_1, \ldots, n + h_k$ are all prime, for any admissible tuple $(h_1, \ldots, h_k)$. Maynard's methods give the correct order of magnitude for the main term but do not yet achieve the conjectured error term. Improving the error term would have consequences for the distribution of prime gaps and for the analysis of prime constellations.

Recent work by Maynard on the \textbf{distribution of primes in function fields} provides a model for the number field case. Over function fields $\mathbb{F}_q[t]$, the analogue of the twin prime conjecture can be studied using geometric methods, and Maynard's sieve methods translate directly. The function field setting provides a testing ground for new ideas that may eventually apply to the integer case.

\begin{center}\rule{0.5\linewidth}{0.5pt}\end{center}

\section*{Summary}
\addcontentsline{toc}{section}{Summary}

Maynard proved that there are infinitely many bounded gaps between primes using a multidimensional sieve method, improved the GPY method, and established the best known lower bounds for large gaps between primes. His work transformed analytic number theory and provided new tools for studying the fine structure of the distribution of primes in arithmetic progressions.
